\section{The Bell interpretive commitment}
\label{sec:bell-interpretation}

The substrate of this paper is deterministic (the equations of motion~\eqref{eq:eom} are second-order
hyperbolic PDEs) and local (all forces propagate through nearest-neighbor spring contacts at finite
speed $c_T$ or $c_L$).
Bell's theorem then implies that the joint distribution of hidden variables and detector settings
cannot factorize: measurement independence must be relaxed.
Superdeterminism is excluded by the substrate's lack of any mechanism to coordinate microscopic
initial conditions with macroscopic measurement choices at spacelike separation.
The unique remaining option consistent with this ontology is \emph{retrocausality}: future measurement
contexts propagate backward along soliton worldtubes to a shared V-vertex in the past, without any
spacelike signal.

The brane action is time-symmetric: its Lorentzian sign structure selects a kind of direction
(timelike vs spacelike) but not a direction along the timelike axis.
The arrow of causality is therefore a property of soliton solutions, not of the substrate.
Matter solitons propagate forward; antimatter solitons propagate backward, realizing the
Feynman--Stueckelberg interpretation at the soliton level.
Entanglement is the continuity of a single branching 4D world-volume that splits at a V-vertex in
the shared past of the entangled pair.

This retrocausal worldtube stance is imported by all companion papers as a fixed ontological
commitment.
A full development of the Bell constraint, the uniqueness argument, and the worldtube
interpretation is given in Paper~V~\citep{MolzbergerBell2026}.